\documentclass{article}
\begin{document}
\section{Hypothesis}
Will the addition of mechanics that allow the player to traverse the map in different ways
allow the use of bigger maps or will it cause the player to be frustrated with the map despite
the extra mechanics to traverse the map.



\section{Goals}
General Goal List:
\begin{itemize}
    \item Connect all the prototypes.
    \item Add multiple ways to traverse the procedurally generated map.
\end{itemize}
\section{Level Design}
\subsection{Story}
The player will find themselves in a sort of laboratory where a disembodied voice will try 
to guide the player out of the laboratory. If the player follows the voice they will successfully
escape the laboratory. If the player disobeys the voice they will be taken deeper into the laboratory
where they will then have to escape from the laboratory without the help of the narrator. 
\subsection{Game Progression}
The general game progression is that the player will start in the tutorial and then make their ways
from the first level to the third level. The first level will be handmade and will have
a narrator. The second level will be a procedurally generated level.
\subsection{Tutorial}
The player will start in the tutorial where they will learn the basics of the game. The player
will first learn the basic controls which are: Moving, sprinting, toggling the player light and
general interactions the player can perform. The player will be introduced to the enemy in the
tutorial and different ways to avoid them.
\subsection{First Level}
The player will be put in a laboratory map where a disembodied voice called the narrator will then
try to help the player escape from the laboratory. If the player follows the narrator then they will
escape from the laboratory, if the player disobeys the narrator they will go deeper into the laboratory.
The enemy will not appear in this level. This level will not be procedurally generated.
\subsection{Second Level}
The player will be in a laboratory that is decayed. The player will face an enemy that is constantly trying
to catch them. The player will also need to try and solve some sort of puzzle which could be finding a code 
or bringing an object from point A to point B. The decay is shown by giving the procedural generation a high level of
random walk.

\section{Genre}
The genre of the game will be a stealth survival horror and puzzle genre. 
The player will need to hide and evade an unstoppable enemy while trying to complete some sort of goal.
This will be similar to the Alien from Alien Isolation, because the AI enemy was inspired from Alien 
Isolation.
The puzzle genre appears when the player will need to find solutions to the puzzles presented to them.

\section{Possible Puzzles}
These are the puzzles that can appear during the procedural generation of the level. Each puzzle will 
have success criteria which is determined by the puzzle.
\subsection{Fetch quest}
The player will need to find an item and bring it to the destination.
\newline
The generation will be a success if:
\begin{itemize} 
    \item The item can be accessed.
    \item The item can be picked up and put down.
    \item The destination can be accessed.
    \item The destination recognizes the item.
    \item The player completes the goal after brining the item to the destination.
\end{itemize}
\subsection{Password}
The player will need to find a password that they will then use to open a lock.
\newline
The generation will be a success if:
\begin{itemize} 
    \item All parts of the password can be accessed
    \item All parts of the password can be read.
    \item The lock can be accessed.
    \item The lock successfully recognizes the entered password.
    \item The lock completes the goal after the password has been entered.
\end{itemize}


\section{Prototype Improvements}
This prototype will improve upon all previous prototypes besides prototype 3, while connecting all the prototypes together.
\subsection{Prototype 1 Improvements}
The map of prototype 1 will be improved while also allowing the player to use more of the mechanics to
explore the map. This will include areas that can only be accessed when the player's light is turned off
and areas that have a timer that the player needs to sprint towards to access. This will give the player
more interaction in exploring the map.
\subsection{Prototype 2 Improvements}
The map of prototype 2 will be improved to make the tutorial a more optional experience instead of a forced
experience. It will gradually show the player the mechanics without overwhelming them by showing them all
the mechanics at once. A small tutorial will be shown to the player in the pause ui so the player can
be reminded about the controls at any point during the game.
\subsection{Prototype 4 Improvements}
The procedural generation algorithm will be improved by keeping track of all the rooms in the generation
to allow adding of special rooms which will either be a room in which a puzzle needs to be solved or 
will have an interaction for the player. The algorithm will also allow for secret passages between
rooms. 

\section{Communication Design}
\subsection{Prototype 4}
To indicate that the enemy knows about the player's position a radar ping sound (reference sound) plays every
few seconds when using the minimap.
Each intractable is highlighted with a green light (change colours to be the same) this shows the player that
they can interact with anything in a green light. A popup also appears prompting the user to press space which
also helps the player know they can interact with specific game elements.


\section{Process}
\subsection{Tutorial Design - Prototype 2}
The design philosophy for the tutorial is that it needs to introduce the player to the general mechanics.
This means it first needs to show the player how to move and then after confirming that the player knows how
to move will then allow the player to select which mechanic they want to learn about. This allows the player
freedom in the tutorial 
\subsection{Prototype 1}
No timer areas were added to the level because the first level is made to feel like the player is calmly being
guided through the level while the timer areas will just add unnecessary anxiety to the player.
\subsection{Prototype 3}
Fixed a bug in which caused the enemy to move as if idle when being sent to specific spots on the map. This bug
sometimes stopped the enemy from moving towards the player when using the minimap causing the minimap to be less
risky to use.
The map made for prototype 3 was just to show off the enemy thus it has been removed making the final level
progression of the game be: Prototype 2, Prototype 1, Prototype 4.

\subsection{Prototype 4}
Increased the player light size because the original light size made the player feel too cramped in the room
causing the player to get lost in the map a bit too often. The original size was a 4 unit circle while the new
size is a 7 unit circle. The detection for the enemy to see the player light is still the original size to allow
the player some leeway in not getting spotted by the enemy but still being able to see what is happening.
A minimap was added to the game to help to give the player a sense of direction, because players got confused
with the maze-like map structure. The minimap will ping the players position every few seconds which will guide 
the enemy to the player. This is to make the minimap a risk for the player to use for too long while still helping
the player navigate the level.


Text was added to the puzzle solution to help the player remember which parts of the solution the player has already found.
This was done because the players would usually forget the solution parts they have found because the enemy chasing them
causes the players to forget what they just found in order to try escape the enemy.

\section{Hypothesis reflection}
\subsection{Mechanics as story}
During the first level the player is guided through the laboratory. There are these black spots that can sometimes
block the player's path where the player will then need to find a way around them. This is where the player will use
some of the movement mechanics added. These include turning off the player light to find secret paths and going to 
teleporters to teleport around the obstacle. Here the movement mechanics must be used to progress instead of just
to traverse the map or run away. These mechanics show the player that the map can be traversed in different ways,but
the player needs to use these traversals to continue the story.
These mechanics in this situation allows for a bigger map because of the narrator which is guiding the player making
the player not feel frustrated as they traverse the map, because they are being led through the map and are not 
in any danger. This can feel boring to players due to the games minimal graphics, but if the graphics were improved
the boring areas can be turned into areas where environmental storytelling keeps the players attention.

\subsection{Mechanics in procedural generation}
The main room connections for the procedural generation algorithm only allows for the rooms to connect in one direction
causing rooms to typically only have one entrance and exit.This causes some of the adjacent rooms to not be connected 
to each other. The added movement mechanics of secret paths allow the player to go to these adjacent rooms. 
The secret paths shorten the time a player would take to get from one room to another while allowing the player to 
actively feel like they are creating these shortcuts. Due to the shortcuts being made from the player toggling their
light it does allow for an escape from the enemy when the player turns off their light to hide from the enemy.
The teleporters also teleport the player from one side of the map to the opposite side of the map. They help the
player traverse the map but also help the player escape from the enemy. The only risk they have is that you don't 
know what is on the other side of the teleporter which can cause the player to instantly be caught by the enemy.

These mechanics in this situation only allow for a bigger map if they are placed in the procedural generation
frequently but not too frequent, this is because of the enemy chasing the player. The mechanics would allow for the player
to be able to escape from the chase but will still need the player to try to hide from the enemy instead of relying on
the mechanics. This will keep the tense atmosphere of the enemy chasing the player while allowing for the player to escape
if they are prepared for the enemy. In this situation the minimum graphics does not take away from the experience, because
the player is already constantly busy with an objective and is constantly wary of the enemy. A graphics improvement will
help with the horror themes of the game.

\section{Reflections}
\subsection{Prototype 1}
In the original prototype 1 the main focus was on the dialog system where the dialog system would try to keep the player
invested in the game by telling a story that the player will be curious about. This prototype only allowed the player to
move which was the only mechanic. Where the player moved was decided by if the player followed the narrator or not.
The improved prototype 1 still keeps the dialog system that tells the player a story and guides the player through the
level where the player can decide if they want to listen to the dialog system or not. 
Some improvements that were made upon the prototype are the inclusion of the mechanics in the game. 
These include toggling the player light to get paths certain areas or using the teleporter to progress, 
while also needing to pick up and item and move it to another area. These inclusion of mechanics help the player learn 
the mechanics and make the level more interesting instead of listening or not listening to the dialog system, 
because the player is actively moving through the level using other mechanics instead of just moving.
The level was also changed from being rigid in areas you are supposed to be and random in areas you are not supposed to
be to rigid everywhere. Black blocks where added to indicate where you should not go. These map changes halp the player
understand more from where they should and shouldn't be according to the story while not confusing the player about the 
map design. (REDO)

\subsection{Prototype 2}
The original prototype 2 lead the player through the tutorial by introducing the mechanics one by one forcing the player
to learn the first mechanic before being able to learn more mechanics. This was really restrictive, thus the newer prototype
tutorial allows the player to learn any mechanic whenever they want after they learnt how to move around. This allows the 
player to quickly look at a mechanic that they want to know more about instead of having to go through the entire tutorial
again. This improvement gives the player more freedom in what they want to learn about the game and when they want to learn
about it. 
\subsection{Prototype 3}
There was a bug with the enemy when the director or minimap tried to move the enemy closer or further away from the 
player, but the idle enemy state overwrote that action and moved the enemy around itself.
This bug did sometimes make the enemy feel as if it was hunting the player when the player was hiding in a hiding spot.
Causing the enemy to look like they are looking around the hiding spot for the player. When the enemy did this it was 
inconsistent, but it needed to be changed because of the addition of the minimap mechanic. The enemy did not 
constantly move towards the player when the minimap was active, because of this bug the minimap did not have the
intended risk of use where the player finds out more about their surrounding but causes the enemy to go to their
current position. Thus the bug had to be removed.
The bug could be turned into a feature, but only when the enemy has lost sight of the player or when searching
around a hiding spot where the player may be. This feature will improve the enemy and make it feel like its 
hunting the player.
\subsection{Prototype 4}
The original prototype 4 was only the procedural generation with an enemy and a fetch quest puzzle. The improved
prototype added more mechanics to the procedural generation allowing the player to find hidden paths to move around
and allowing the player to easily get to the other side of the map. The improved prototype is more interactive with
the player where the player needs to actively search for the puzzle parts and actively activate the secret paths.
These interaction improved the players enjoyment of the game.
A minimap was added to help the player navigate, but this minimap took away from the game by giving the player a 
sense of direction. This takes away from the horror elements of the game because the player feels a sense of dread
because they are lost, which the minimap diminishes. The current procedural generation generates the map in a maze-like
way which may be too confusing for players, thus the minimap helps the players traverse the map. Both improvements want
to be kept thus a solution was found. The solution is that if the player uses the minimap the enemy will be alerted to
their position and will constantly go to the position of the player. This solution makes the minimap a limited resource
that the player can use at a risk. This leads to the horror elements of the game staying by letting the player feel like
they are in danger because of the use of the minimap, but it also helps the player traverse the maze-like map reducing
the player's frustration.


\section{references (delete subsections later)}
\subsection{Tutorials}

https://medium.com/@doandaniel/gamedev-protips-a-few-helpful-tips-on-crafting-a-better-tutorial-ffcc636039d3
https://www.gamedeveloper.com/business/gamedev-thoughts-how-to-design-a-good-tutorial-for-your-indie-game
\subsection{Minimap}


https://pixabay.com/sound-effects/film-special-effects-radar-ping-306440/

\end{document}